% 图 65.2 —— 由 `checks/emit_hand.py` 自动拼出（probe 精测 + prims2 图元分解）
%%
% ⚠ 这是**稿子不是成品**：跑 `checks/checkhand.py 65.2` 看指标 + 看叠合图。
%%
%% ⚠ 待核：
%   · 本图 probe 报出阴影族 1 个 —— **本脚本不生成阴影**（要照原书手写）；prims2 的 strokes 里可能含阴影线，核对时留意别画重
%   · 可疑字形 1 项（空心圈 ⊙ 常被认成字母 o）—— 见文件末
%%
\documentclass[12pt]{standalone}
\usepackage{fontspec}
\setsansfont{Arial}
\newfontfamily\mfallback{DejaVu Sans}
\usepackage{tikz}
\begin{document}
\begin{tikzpicture}[x=1pt, y=-1pt, line width=5.0pt, line cap=round, line join=round]

  %% ════ 直线段（prims2 的 strokes）════
  \draw[line width=5.0pt] (161.0,220.0) -- (220.0,220.0);
  \draw[line width=4.0pt] (229.0,164.0) -- (229.0,211.0);
  \draw[line width=4.0pt] (656.0,219.0) -- (747.0,220.0);
  \draw[line width=5.0pt] (230.0,82.0) -- (231.0,21.0);
  \draw[line width=5.0pt] (238.0,217.0) -- (230.0,212.0);
  \draw[line width=5.0pt] (230.0,86.0) -- (230.0,160.0);
  \draw[line width=4.3pt] (867.0,339.0) -- (866.0,335.0);
  \draw[line width=4.0pt] (998.5,183.0) -- (1010.1,181.9);
  \draw[line width=4.0pt] (934.9,85.0) -- (914.7,92.0);
  \draw[line width=5.0pt] (914.7,92.0) -- (869.0,84.0);
  \draw[line width=4.0pt] (868.0,85.0) -- (867.0,210.0);
  \draw[line width=5.0pt] (228.0,212.0) -- (220.0,220.0);
  \draw[line width=5.0pt] (867.0,341.0) -- (867.0,386.0);
  \draw[line width=4.0pt] (220.0,220.0) -- (229.0,230.0);
  \draw[line width=4.0pt] (766.2,81.2) -- (768.0,69.0);
  \draw[line width=4.0pt] (851.0,74.0) -- (866.0,79.0);
  \draw[line width=4.0pt] (228.0,280.0) -- (215.4,281.0);
  \draw[line width=4.0pt] (230.0,329.0) -- (229.0,281.0);
  \draw[line width=5.0pt] (442.0,220.0) -- (354.0,220.0);
  \draw[line width=5.0pt] (352.0,221.0) -- (283.0,221.0);
  \draw[line width=3.0pt] (1094.0,84.0) -- (1074.0,84.0);
  \draw[line width=5.0pt] (866.0,334.0) -- (867.0,229.0);
  \draw[line width=4.0pt] (866.0,17.0) -- (867.0,78.0);
  \draw[line width=3.3pt] (867.0,80.0) -- (868.0,83.0);
  \draw[line width=4.0pt] (939.0,219.0) -- (876.0,219.0);
  \draw[line width=4.0pt] (229.0,231.0) -- (229.0,278.0);
  \draw[line width=4.5pt] (866.0,211.0) -- (858.0,219.0);
  \draw[line width=4.5pt] (278.0,219.0) -- (273.7,206.7);
  \draw[line width=4.0pt] (273.7,206.7) -- (267.0,196.5);
  \draw[line width=5.0pt] (347.0,180.9) -- (352.0,198.8);
  \draw[line width=4.0pt] (352.0,198.8) -- (354.0,220.0);
  \draw[line width=5.0pt] (159.0,221.0) -- (106.0,221.0);
  \draw[line width=4.0pt] (943.0,221.0) -- (946.0,239.8);
  \draw[line width=5.0pt] (873.0,340.0) -- (868.0,340.0);
  \draw[line width=4.5pt] (866.0,228.0) -- (858.0,219.0);
  \draw[line width=5.0pt] (868.0,211.0) -- (876.0,219.0);
  \draw[line width=5.0pt] (282.0,230.0) -- (282.0,222.0);
  \draw[line width=4.8pt] (230.0,230.0) -- (234.2,227.8);
  \draw[line width=5.0pt] (234.2,227.8) -- (237.8,221.2);
  \draw[line width=5.0pt] (237.8,221.2) -- (278.0,220.0);
  \draw[line width=4.0pt] (858.0,219.0) -- (753.0,218.0);
  \draw[line width=4.7pt] (747.0,220.0) -- (751.0,218.0);
  \draw[line width=5.0pt] (756.1,343.0) -- (761.7,344.0);
  \draw[line width=5.0pt] (876.0,219.0) -- (867.0,228.0);
  \draw[line width=5.0pt] (20.0,220.0) -- (104.0,220.0);
  \draw[line width=4.0pt] (1061.0,219.0) -- (943.0,220.0);
  \draw[line width=5.0pt] (229.0,332.0) -- (229.0,387.0);

  %% ════ 曲线（prims2 的 curves —— 三段式，坐标同为 (y,x)）════
  \draw[line width=4.0pt] (228.0,163.0) .. controls (199.3,169.3) and (162.3,185.5) .. (161.0,220.0);
  \draw[line width=4.0pt] (231.0,83.0) .. controls (263.5,79.6) and (320.4,72.5) .. (322.0,118.2);
  \draw[line width=4.0pt] (322.0,118.2) .. controls (314.4,161.1) and (262.4,151.7) .. (231.0,161.0);
  \draw[line width=2.0pt] (1093.0,77.0) .. controls (1087.1,78.8) and (1079.0,78.5) .. (1073.0,77.0);
  \draw[line width=4.0pt] (940.0,218.0) .. controls (936.4,193.8) and (911.5,168.0) .. (929.4,143.6);
  \draw[line width=5.0pt] (929.4,143.6) .. controls (936.3,136.2) and (947.7,131.8) .. (956.4,138.4);
  \draw[line width=5.0pt] (956.4,138.4) .. controls (973.7,151.1) and (978.8,174.4) .. (998.5,183.0);
  \draw[line width=5.0pt] (1010.1,181.9) .. controls (1037.6,159.5) and (1034.0,114.7) .. (1015.6,87.6);
  \draw[line width=4.0pt] (1015.6,87.6) .. controls (995.3,62.6) and (956.8,64.1) .. (934.9,85.0);
  \draw[line width=4.0pt] (105.0,219.0) .. controls (102.0,150.3) and (161.3,87.9) .. (229.0,85.0);
  \draw[line width=4.0pt] (752.0,217.0) .. controls (763.1,181.4) and (810.9,180.1) .. (827.0,149.9);
  \draw[line width=5.0pt] (827.0,149.9) .. controls (828.7,134.2) and (816.8,122.4) .. (805.3,113.3);
  \draw[line width=4.0pt] (805.3,113.3) .. controls (789.9,106.1) and (769.0,101.0) .. (766.2,81.2);
  \draw[line width=4.0pt] (768.0,69.0) .. controls (788.9,43.7) and (829.3,55.0) .. (851.0,74.0);
  \draw[line width=4.0pt] (215.4,281.0) .. controls (182.1,288.5) and (150.1,255.5) .. (160.0,222.0);
  \draw[line width=5.0pt] (231.0,330.0) .. controls (286.1,323.5) and (352.5,283.8) .. (353.0,222.0);
  \draw[line width=4.0pt] (105.0,222.0) .. controls (106.9,256.3) and (113.5,291.8) .. (142.1,315.0);
  \draw[line width=4.0pt] (142.1,315.0) .. controls (167.1,330.7) and (197.8,334.0) .. (228.0,331.0);
  \draw[line width=4.0pt] (267.0,196.5) .. controls (267.1,160.3) and (334.2,144.9) .. (347.0,180.9);
  \draw[line width=4.0pt] (946.0,239.8) .. controls (952.3,262.6) and (953.9,290.3) .. (944.6,312.4);
  \draw[line width=5.0pt] (944.6,312.4) .. controls (928.9,335.3) and (899.6,353.3) .. (873.0,340.0);
  \draw[line width=5.0pt] (230.0,279.0) .. controls (254.0,275.8) and (288.6,260.1) .. (282.0,230.0);
  \draw[line width=5.0pt] (747.0,220.0) .. controls (727.9,257.1) and (723.0,311.4) .. (756.1,343.0);
  \draw[line width=4.0pt] (761.7,344.0) .. controls (786.4,342.4) and (790.8,312.1) .. (808.4,301.0);
  \draw[line width=4.0pt] (808.4,301.0) .. controls (831.9,297.8) and (843.3,329.8) .. (865.0,334.0);

  %% ════ 标签（probe 直接给的 \node 行）════
  %% ⚠ 全部补了 fill=white：文字在最上层，否则与曲线交叉干扰
     \node[anchor=center,inner sep=0pt,fill=white,font=\fontsize{31.1}{38.8}\selectfont] at (441.3,76.8) {\textsf{\textbf{6}}};   % box(433,65,15,22)
     \node[anchor=center,inner sep=0pt,fill=white,font=\fontsize{30.9}{38.7}\selectfont] at (380.8,80.3) {\textsf{\textbf{\textit{n}}}};   % box(372,70,16,17)
     \node[anchor=center,inner sep=0pt,fill=white,font=\fontsize{29.5}{36.9}\selectfont] at (1137.9,81.5) {\textsf{\textbf{1}}};   % box(1130,70,9,22)
     \node[anchor=center,inner sep=0pt,fill=white,font=\fontsize{30.3}{37.9}\selectfont] at (417.8,78.6) {{\mfallback\textbf{−}}};   % box(402,72,17,4)
     \node[anchor=center,inner sep=0pt,fill=white,font=\fontsize{52.4}{65.5}\selectfont] at (1118.7,85.0) {\textsf{\textbf{|}}};   % box(1109,72,16,19)
     \node[anchor=center,inner sep=0pt,fill=white,font=\fontsize{30.0}{37.5}\selectfont] at (1051.8,84.8) {\textsf{\textbf{\textit{n}}}};   % box(1043,75,16,16)
     \node[anchor=center,inner sep=0pt,fill=white,font=\fontsize{30.3}{37.9}\selectfont] at (417.8,85.6) {{\mfallback\textbf{−}}};   % box(402,79,17,4)

  % 钉死画布：否则超出的图元/控制点会把页面撑大，1:1 回比就失准
  \pgfresetboundingbox
  \path[use as bounding box] (0,0) rectangle (1154,401);
\end{tikzpicture}
\end{document}

% ⚠ 待核清单（probe 拿不准，人眼过一遍）：
%   · 未识别/跳过：⑤ 跳过/未识别的字形
